\documentclass[11pt]{article}
\usepackage[letterpaper,margin=0.6in,top=0.55in,bottom=0.6in]{geometry}
\usepackage[T1]{fontenc}
\usepackage{lmodern,microtype,enumitem,xcolor,fancyhdr}
\renewcommand{\familydefault}{\sfdefault}
\setlist{nosep,leftmargin=1.2em,itemsep=1.5pt}
\setlength{\parindent}{0pt}\setlength{\parskip}{4pt}
\pagestyle{fancy}\fancyhf{}\renewcommand{\headrulewidth}{0pt}
\fancyfoot[L]{\footnotesize Bellingham Math Circle / Week 34 / Adult table card / Unpiloted}
\fancyfoot[R]{\footnotesize Full answers: the facilitator guide}
\newcommand{\sect}[1]{\par\medskip{\large\bfseries #1}\par\smallskip}
\begin{document}
{\footnotesize Week 34 / Hidden turns / For the adult at each table}\par
{\LARGE\bfseries Hidden turns: the one-page version}\par
\sect{What the activity is}
A child fills a printed ring of spots with coloured counters and makes a tracing-paper copy of the pattern. A partner turns or flips the copy to a \emph{new} position and tries to make every dot land on its own colour. The builder's goal is a pattern that can only match when the copy goes straight back where it started, using as few kinds of counter as possible.

\sect{Materials (about 15--20 minutes)}
\begin{itemize}
\item Per pair of children: about 24 counters, 8 each of three kinds that are easy to tell apart (at least 20~mm wide), in a small bowl.
\item Per pair: about 10 sheets of tracing paper big enough to cover a ring, and coloured pencils in the counters' colours.
\item Print the ring pages at actual size. At home, make one dotted copy and check that you can still see the dots and colours when the sheet is flipped over.
\end{itemize}

\sect{Launch (one minute, whole group)}
On a four-spot ring, lay out two kinds alternating (red, blue, red, blue). Press a coloured dot onto tracing paper over each spot. Keep the counters still and move only the copy: turn it one spot (\emph{no match}), turn it two spots (\emph{match}), then flip it over the line through two opposite spots (\emph{match}). Say: ``Putting it back where it started doesn't count. Can you build a pattern your partner can't match?''

\sect{What Problem 1 is for}
\begin{itemize}
\item \textbf{K--1:} a three-spot ring with two kinds. The builder \emph{cannot} win: whichever counter is the odd one out, a flip through it always matches. That is the point. Swap roles after each catch; when a child says ``the flip always gets me,'' agree and go to Problem~2, where a third kind lets the builder win.
\item \textbf{Grades 2--3 and 4--5:} the same three-ring, asking for the fewest kinds. The answer is 3 (one of each); with two kinds a flip always matches. Grades 4--5 also explain why two kinds can never work.
\end{itemize}

\sect{The surprise}
Small rings (3, 4 or 5 spots) need three kinds of counter, but every ring with 6 or more spots can be done with just two. On the bigger rings, three spots of one kind at uneven distances apart (gaps of 1, 2 and the rest) block every turn and flip.

\sect{Key answers (A, B, C are any three kinds, read clockwise)}
\begin{itemize}
\item \textbf{K--1.} P1: impossible with two kinds. P2 (four-ring): 3 kinds, e.g.\ A B C C. P3: first tell children which counter each printed mark stands for (filled dot, ring, square). Two changes keep it unmatched: the top dot to the square's kind, or the right dot to the ring's kind. P4 (five-ring): two kinds fail; A B C C C works. P5 (six-ring): A A B A B B. P6 (eight-ring): A A B A B B B B.
\item \textbf{Grades 2--3.} Three-, four- and five-rings: 3 kinds each (A B C; A B C C; A B C C C). P4: A A B A B B. P5: A A B A B B B B; the smaller kind needs at least 3 counters.
\item \textbf{Grades 4--5.} P1--3: 3 kinds; with two kinds the smaller kind has at most two spots, and some flip always swaps those two spots (or holds a single one still), so it matches. P4: A A B A B B. P5: on any ring of 6 or more spots, put A on spots 0, 1 and 3 and B everywhere else; the gaps 1, 2 and the rest are all different, so no turn or flip matches.
\end{itemize}
\end{document}
